% from ega4-16-fr.tex
\oldpage[IV-4]{22}
identifient $S$ au produit des $(X\times_S X)$-préschémas $X$ et $X$ pour les
morphismes $\Delta_f$ et $(g\circ f,1_X)_S$ (resp.\ $(1_X,g\circ f)_S$).
D'autre part, les diagrammes
\[
  \xymatrix@=3pc{
    X\ar[r]^{(g\circ f,1_X)_S}\ar[d]_-f&X\times_S X\ar[d]^-{p_1}
      &X\ar[r]^{(1_X,g\circ f)_S}\ar[d]_-f&X\times_S X\ar[d]^-{p_2}\\
    S\ar[r]_{g}&X&S\ar[r]_{g}&X
  }
  \tag{16.4.11.4}
  \label{IV.16.4.11.4}
\]
identifient $X$ au produit des $X$-préschémas $S$ et $X\times_S X$ pour les
morphismes $g$ et $p_1$ (resp.\ $p_2$) (cas particulier de la formule
d'associativité (I, 3.3.9.1)). On peut dire que $\Delta_f$, considéré comme
$X$-section de $X\times_S X$ (relativement à $p_1$ ou $p_2$) joue le rôle
d'une \emph{section universelle} pour les $S$-sections de $X$: chacune de ces
sections $g$ s'en déduit en effet par le \emph{changement de base}
$(g\circ f,1_X)_S:X\to X\times_S X$. La définition de l'homomorphisme
$\varpi_n$ et le fait qu'il est bijectif résultent donc de ces remarques et de
(16.2.3, (ii)) appliqué au premier diagramme (16.4.11.4). La commutativité du
diagramme (16.4.11.2) résulte de même de (16.2.3, (ii)) appliqué cette fois au
second diagramme (16.4.11.4). Pour expliciter $\varpi_n$, on peut se borner au
cas où $g$ est une immersion \emph{fermée}: en effet, pour tout $s\in S$, il y
a un voisinage ouvert $W$ de $s$ dans $S$ tel que $g(W)$ soit fermé dans un
ouvert $U$ de $X$, et il est clair que $g|W$ est une $W$-section du morphisme
$U\cap f^{-1}(W)\to W$, restriction de $f$, et $g(W)$ est \emph{a fortiori}
fermé dans $U\cap f^{-1}(W)$. On peut donc supposer que $S$ est un
sous-préschéma fermé de $X$ défini par un Idéal quasi-cohérent $\mathscr{K}$.
Alors les définitions précédentes montrent que si $W$ est un ouvert de $S$,
$t$ une section de $\mathcal{O}_X$ au-dessus de $f^{-1}(W)$,
$\varpi_n(d^n t|W)$ est égal à l'image canonique de $t$ dans
$\Gamma(W,(\mathcal{O}_X/\mathscr{K}^{n+1})|W)$. L'unicité de $\varpi_n$
résulte donc de ce que l'image de $\mathcal{O}_X$ par $d_{X/S}^n$ engendre le
$\mathcal{O}_X$-Module $\mathscr{P}_{X/S}^n$ (16.3.8).

\begin{corollary}[16.4.12]
\label{IV.16.4.12}
Soient $k$ un corps, $X$ un $k$-préschéma, $x$ un point de $X$ rationnel sur
$k$. Alors $(\mathscr{P}_{X/k}^n)_x\otimes_{\mathcal{O}_x}k(x)$ est
canoniquement isomorphe (en tant que $k(x)$-algèbre augmentée) à
$\mathcal{O}_x/\mathfrak{m}_x^{n+1}$.
\end{corollary}

Il suffit de considérer l'unique $k$-section $g$ de $X$ telle que
$g(\operatorname{Spec}(k))=\{x\}$.

\begin{corollary}[16.4.13]
\label{IV.16.4.13}
Soient $f:X\to S$ un morphisme, $s$ un point de $S$,
$X_s=X\times_S\operatorname{Spec}(k(s))$ la fibre de $f$ en $s$. Si
$x\in X_s$ est rationnel sur $k(s)$,
$(\mathscr{P}_{X/S}^n)_x\otimes_{\mathcal{O}_s}k(s)$ est canoniquement
isomorphe à $\mathcal{O}_{X_s,x}/\mathfrak{m}'_x{}^{\,n+1}$, où
$\mathfrak{m}'_x$ est l'idéal maximal de $\mathcal{O}_{X_s,x}$; de façon
précise, cet isomorphisme fait correspondre à $(d^n t)_x\otimes1$ (où $t$ est
une section de $\mathcal{O}_X$ au-dessus d'un voisinage ouvert de $x$ dans
$X$) la classe de $t_x\otimes1$ modulo $\mathfrak{m}'_x{}^{\,n+1}$.
\end{corollary}

Cela résulte de (16.4.5) et (16.4.12).

Les corollaires précédents justifient la terminologie de «faisceau des parties
principales d'ordre $n$».

\begin{proposition}[16.4.14]
\label{IV.16.4.14}
Soit $\rho:A\to B$ un homomorphisme d'anneaux, $S$ une partie
multiplicativement stable de $B$. Alors les homomorphismes canoniques
\[
  S^{-1}P_{B/A}^n\longrightarrow P_{S^{-1}B/A}^n
  \tag{16.4.14.1}
  \label{IV.16.4.14.1}
\]
